\clearpage
\section{Phân tích \& Thiết kế hệ thống}
\subsection{Kiến trúc tổng thể hệ thống}

Hệ thống UTraffic được thiết kế theo mô hình \textbf{microservice} kết hợp với \textbf{kiến trúc Kappa} (Kappa Architecture) nhằm đảm bảo khả năng xử lý dữ liệu giao thông thời gian thực, đồng thời hỗ trợ các tác vụ phân tích và dự báo phức tạp dựa trên mô hình học sâu trên đồ thị (T-GCN). Toàn bộ hệ thống được tổ chức thành 5 tầng chức năng độc lập nhưng liên kết chặt chẽ với nhau, bao gồm: tầng thu thập dữ liệu, tầng xử lý luồng, tầng lưu trữ lai ghép, tầng học máy không gian - thời gian, và tầng phục vụ.

\begin{figure}[H]
    \centering
    \includegraphics[width=\textwidth]{images/architecture.png}
    \caption{Kiến trúc tổng thể hệ thống UTraffic}
    \label{fig:architecture}
\end{figure}

\subsection{Tầng thu thập dữ liệu (Traffic Data Sources)}

Hệ thống tiếp nhận dữ liệu giao thông từ hai nguồn chính:

\begin{itemize}
    \item \textbf{TomTom Traffic Stats API:} Cung cấp dữ liệu tốc độ thực tế theo từng phân đoạn đường (road segment) với độ phân giải thời gian 15 phút. Dữ liệu bao gồm các thuộc tính: tốc độ trung bình (\texttt{average\_speed}), tốc độ điều hòa (\texttt{harmonic\_average\_speed}), tốc độ trung vị (\texttt{median\_speed}), thời gian di chuyển (\texttt{average\_travel\_time}), tỉ lệ thời gian di chuyển so với điều kiện lý tưởng (\texttt{travel\_time\_ratio}), và kích thước mẫu (\texttt{sample\_size}). Dữ liệu được thu thập theo mô hình batch job từ TomTom Movement Stats API, bao gồm 31 ngày × 24 khung giờ 15 phút/ngày cho các phân đoạn đường thuộc khu vực nghiên cứu tại TP.HCM (bounding box: lat 10.760--10.805, lon 106.675--106.715).

    \item \textbf{OpenStreetMap API (OSMnx):} Cung cấp dữ liệu topo học đường bộ dưới dạng đồ thị định hướng, bao gồm các node (giao lộ) với tọa độ địa lý và các edge (đoạn đường) với các thuộc tính như chiều dài, tốc độ giới hạn, số làn xe, và loại đường (\texttt{highway type}). Dữ liệu OSM được lưu cache dưới định dạng \texttt{.graphml} để tránh tải lại mỗi lần chạy pipeline.
\end{itemize}

\subsection{Tầng xử lý luồng (Kappa Architecture -- Stream Processing Engine)}

Kiến trúc Kappa được lựa chọn thay cho kiến trúc Lambda nhằm đơn giản hóa hệ thống bằng cách xử lý tất cả dữ liệu thông qua một luồng duy nhất, loại bỏ sự phức tạp của việc duy trì song song hai lớp batch và streaming. Tầng xử lý bao gồm các thành phần sau:

\subsubsection{Data Ingestion Layer}
Là điểm tiếp nhận dữ liệu từ các nguồn bên ngoài vào hệ thống. Thành phần này chịu trách nhiệm xác thực định dạng đầu vào và phân phối dữ liệu thô vào hàng đợi xử lý. Trong giai đoạn triển khai hiện tại, Data Ingestion Layer hoạt động theo mô hình kéo dữ liệu định kỳ (scheduled pull) từ TomTom API thay vì streaming thời gian thực do giới hạn của API phía nhà cung cấp.

\subsubsection{Stream Consumer \& Data Validation}
Stream Consumer tiêu thụ dữ liệu từ hàng đợi và chuyển tiếp sang module Data Validation \& Cleaning. Tại đây, dữ liệu được kiểm tra tính toàn vẹn (null check, range check cho các giá trị tốc độ), loại bỏ các bản ghi có \texttt{sample\_size = 0} (không có xe đi qua), và chuẩn hóa định dạng timestamp theo múi giờ \texttt{Asia/Ho\_Chi\_Minh}.

\subsubsection{Road Topology Builder (OSM) \& Map Matcher}
Module Road Topology Builder sử dụng thư viện \texttt{OSMnx} để xây dựng đồ thị đường bộ từ OpenStreetMap. Đồ thị này sau đó được lưu cache dưới dạng \texttt{.npz} để phục vụ các bước xử lý tiếp theo.

Map Matcher là module trọng tâm của pipeline, thực hiện gán (matching) các phân đoạn TomTom vào các cạnh OSM tương ứng thông qua thuật toán \textbf{shortest-path map matching}. Với mỗi phân đoạn TomTom, hệ thống xác định node OSM gần nhất với điểm đầu và điểm cuối của phân đoạn, sau đó tính đường đi ngắn nhất giữa hai node đó. Tất cả các cạnh OSM trên đường đi này đều nhận thuộc tính giao thông của phân đoạn TomTom tương ứng. Thiết kế này giải quyết vấn đề không khớp tỉ lệ giữa phân đoạn TomTom (thường dài 200--500m) và cạnh OSM (thường dài 30--100m). Kết quả là một \textbf{subgraph} chỉ bao gồm các cạnh OSM đã được xác nhận có dữ liệu TomTom tương ứng, đảm bảo 100\% dữ liệu là thực tế (không có giá trị mặc định giả tạo).

\subsubsection{Feature Engineering}
Module này tính toán các đặc trưng bổ sung từ dữ liệu thô:
\begin{itemize}
    \item \textbf{Congestion Index:} $CI = 1 - \frac{v_{observed}}{v_{limit}}$, với giá trị từ 0 (thông thoáng) đến 1 (tắc hoàn toàn).
    \item \textbf{Speed Limit Ratio:} $SLR = \frac{v_{observed}}{v_{limit}}$, biểu diễn mức độ sử dụng năng lực đường.
    \item \textbf{Std Speed:} Độ lệch chuẩn tốc độ trong khung giờ, phản ánh mức độ ổn định của dòng chảy giao thông.
    \item \textbf{Node-level aggregation:} Đặc trưng topo học của mỗi node (bậc đỉnh, betweenness centrality, tọa độ chuẩn hóa) được tính từ subgraph sau map matching.
\end{itemize}

\subsubsection{Binary .npz Formatter \& Time-Series Sync}
Dữ liệu sau xử lý được lưu dưới hai định dạng song song:
\begin{itemize}
    \item \textbf{Binary .npz:} Lưu trữ tensor đặc trưng dạng \texttt{(T, N, F)} cho T-GCN, trong đó $T$ là số bước thời gian, $N = 385$ là số node trong subgraph, và $F$ là số đặc trưng (9 đặc trưng TomTom + 4 đặc trưng topo học OSM). Định dạng này được tối ưu cho việc nạp dữ liệu vào mô hình học máy.
    \item \textbf{Time-Series Sync:} Ghi dữ liệu quan sát theo thời gian vào TimescaleDB thông qua batch insert, phục vụ truy vấn lịch sử và tính toán aggregate.
\end{itemize}

\subsection{Tầng lưu trữ lai ghép (Hybrid Storage System)}

Hệ thống lưu trữ được thiết kế theo mô hình lai ghép nhằm cân bằng giữa hiệu năng truy xuất, chi phí lưu trữ, và khả năng mở rộng:

\subsubsection{Object Store -- AWS S3 (Data Lake)}
Lưu trữ toàn bộ dữ liệu thô và đã xử lý dưới dạng file nhị phân:
\begin{itemize}
    \item \textbf{Raw Archive (.json/.csv):} Dữ liệu gốc từ TomTom API, lưu nguyên bản để phục vụ tái xử lý khi cần.
    \item \textbf{Map \& Topology (.graphml/.npz):} Đồ thị OSM và kết quả map matching.
    \item \textbf{Processed Features (.npz):} Tensor đặc trưng đã xử lý, sẵn sàng cho huấn luyện mô hình.
    \item \textbf{Model Registry (.pth):} Checkpoint của các mô hình T-GCN và DTC-STGCN đã huấn luyện.
    \item \textbf{Analytics Results (.csv/.png/.json):} Kết quả phân tích tương quan, heatmap, và báo cáo so sánh mô hình.
\end{itemize}

\subsubsection{Fast Access \& DB Layer}
\begin{itemize}
    \item \textbf{Local Disk Cache:} Lưu cache các file \texttt{.npz} thường xuyên truy cập (graph structure, model weights, correlation matrix) trên ổ SSD cục bộ để giảm độ trễ khi inference.
    \item \textbf{TimescaleDB -- Time-Series Traffic:} Mở rộng của PostgreSQL được tối ưu cho dữ liệu chuỗi thời gian. Bảng \texttt{traffic\_observations} được cấu hình là hypertable với chunk theo tuần, cho phép truy vấn dữ liệu lịch sử hiệu quả thông qua continuous aggregates.
    \item \textbf{PostgreSQL -- User \& Nodes \& Metadata:} Lưu trữ dữ liệu có cấu trúc bền vững: thông tin người dùng, metadata node/edge của đồ thị đường bộ, kết quả tương quan, và phiên bản mô hình. Tích hợp extension \textbf{PostGIS} để hỗ trợ truy vấn không gian địa lý.
\end{itemize}

\subsubsection{Redis Cache Layer}
Bốn instance Redis chuyên biệt phục vụ các mục đích khác nhau:
\begin{itemize}
    \item \textbf{Redis-Auth (Token Store):} Lưu JWT blacklist và session metadata. TTL được đặt bằng thời gian hết hạn của token.
    \item \textbf{Redis-Corr (Correlation Cache):} Cache kết quả tương quan top-20 láng giềng của từng node dưới dạng JSON blob. Kích thước cache ước tính $\sim$0.63 MB cho toàn bộ 385 node. TTL được cập nhật theo batch interval khi có correlation job mới.
    \item \textbf{Redis-Pred (Prediction Cache):} Cache dự báo tắc nghẽn cho 4 horizon tiếp theo của từng node. TTL 30 giây để đảm bảo dữ liệu luôn gần với thời gian thực.
    \item \textbf{Redis-API (API Cache):} Cache response của các API endpoint thường xuyên được gọi (danh sách node, thông tin đồ thị). TTL 60 giây.
\end{itemize}

\subsection{Tầng học máy không gian -- thời gian (Spatio-Temporal Learning Engine)}

\subsubsection{Mô hình T-GCN}
Mô hình \textbf{Temporal Graph Convolutional Network (T-GCN)} là thành phần cốt lõi của tầng ML. Kiến trúc T-GCN kết hợp \textbf{Graph Convolutional Network (GCN)} để học mối phụ thuộc không gian giữa các node với \textbf{Gated Recurrent Unit (GRU)} để mô hình hóa sự phụ thuộc thời gian trong chuỗi tốc độ. Cụ thể:

\begin{itemize}
    \item \textbf{Encoder:} Xử lý chuỗi lịch sử $seq\_len = 12$ bước thời gian (tương ứng 3 giờ với bước 15 phút). Tại mỗi bước, TGCNCell thực hiện: (1) GCN tổng hợp thông tin từ các node láng giềng theo ma trận kề chuẩn hóa $\hat{A} = D^{-1/2}AD^{-1/2}$; (2) GRUCell cập nhật trạng thái ẩn; (3) LayerNorm ổn định gradient.
    \item \textbf{Decoder:} Sinh dự báo tự hồi quy (auto-regressive) cho $pred\_len = 4$ bước tiếp theo (tương ứng 60 phút), sử dụng giá trị dự báo của bước trước làm đầu vào cho bước tiếp theo.
    \item \textbf{Output:} Tensor dự báo tốc độ có shape \texttt{(batch, 4, 385, 1)}, trong đó 4 horizon tương ứng với các mốc dự báo T+15, T+30, T+45, T+60 phút.
\end{itemize}

Sau quá trình thực nghiệm và so sánh với các baseline (Historical Average, GRU, SVR, DTC-STGCN), mô hình T-GCN đạt kết quả tốt nhất với RMSE = 0.443 và R\textsuperscript{2} = 0.797 tại horizon 1 bước (15 phút).

\subsubsection{Correlation Jobs}
Module phân tích tương quan sử dụng class \texttt{TGCNCorrelationExtractor} để trích xuất ma trận tương quan nút--nút từ mô hình T-GCN đã huấn luyện. Phương pháp \textbf{detrended Pearson correlation} được áp dụng trên các chuỗi dự báo tốc độ (speed predictions) thay vì trên hidden states, nhằm tránh hiện tượng tương quan giả (spurious correlation) do LayerNorm gây ra. Kết quả là ma trận $385 \times 385$ với giá trị trong khoảng $[-1, 1]$, được lưu vào PostgreSQL (các cặp có $|corr| \geq 0.7$) và file \texttt{.npz} (ma trận đầy đủ).

\subsubsection{Prediction Jobs}
Chạy định kỳ mỗi 15 phút để sinh dự báo cho 4 horizon tiếp theo. Kết quả dự báo được ghi đồng thời vào PostgreSQL (bảng \texttt{predictions}) và Redis-Pred cache để phục vụ API.

\subsubsection{MLflow Experiment Tracking}
Toàn bộ quá trình huấn luyện và đánh giá mô hình được theo dõi bằng MLflow, bao gồm: hyperparameter, metrics theo từng epoch, checkpoint path, và kết quả so sánh giữa các mô hình.

\subsection{Tầng phục vụ (Serving Layer)}

\subsubsection{Backend -- FastAPI}
Backend được xây dựng bằng \textbf{FastAPI} theo mô hình microservice với 4 service độc lập:

\begin{itemize}
    \item \textbf{Auth Service -- JWT:} Quản lý đăng ký, đăng nhập, làm mới token (refresh token rotation), và thu hồi token. Sử dụng Redis-Auth để blacklist JWT khi người dùng đăng xuất.
    \item \textbf{Correlation Service:} Cung cấp API trả về top-K node tương quan với một node cho trước. Ưu tiên đọc từ Redis-Corr cache; trong trường hợp cache miss mới truy vấn PostgreSQL (\texttt{node\_correlation\_cache}).
    \item \textbf{Prediction Service:} Cung cấp API trả về dự báo tắc nghẽn theo horizon cho một hoặc nhiều node. Kết hợp dữ liệu từ Redis-Pred và PostgreSQL.
    \item \textbf{Traffic Service:} Cung cấp API phục vụ bản đồ Folium, bao gồm: truy vấn node/edge theo bounding box (sử dụng PostGIS), thống kê tốc độ theo giờ, và dữ liệu topo học đồ thị.
\end{itemize}

Tất cả các service đều đi qua \textbf{API Gateway} để xử lý xác thực, giới hạn tốc độ (rate limiting), và định tuyến request.

\subsubsection{Nginx / Reverse Proxy}
Nginx đóng vai trò reverse proxy và load balancer, xử lý SSL termination và rate limiting ở tầng mạng. Cấu hình gzip compression cho response JSON nhằm giảm băng thông khi trả về dữ liệu bản đồ.

\subsubsection{Frontend -- ReactJS (UTraffic Platform)}
Giao diện người dùng được xây dựng bằng \textbf{ReactJS} với 4 màn hình chính:

\begin{itemize}
    \item \textbf{Login:} Xác thực người dùng với JWT, hỗ trợ remember me thông qua refresh token.
    \item \textbf{Map View:} Hiển thị bản đồ giao thông tương tác sử dụng \textbf{Folium} (Python) để sinh tile GeoJSON, kết hợp với Leaflet.js trên frontend. Áp dụng kỹ thuật \textbf{lazy loading theo bounding box}: chỉ tải dữ liệu node/edge trong viewport hiện tại, tự động cập nhật khi người dùng pan hoặc zoom bản đồ.
    \item \textbf{Node Detail:} Khi người dùng click vào một node giao thông, hệ thống hiển thị overlay màu sắc biểu diễn mức độ tương quan của node đó với tất cả các node còn lại (đỏ = tương quan cao, xanh = thấp), cùng với biểu đồ dự báo tốc độ 4 horizon tiếp theo.
    \item \textbf{Dashboard:} Tổng quan hệ thống với các thống kê tổng hợp: phân bố tốc độ theo giờ, heat map tắc nghẽn theo thời gian, và danh sách top node có mức độ tương quan cao nhất.
\end{itemize}

\subsection{Tầng giám sát -- DevOps (Monitoring \& Experimentation)}

\begin{itemize}
    \item \textbf{Prometheus Metrics Collector:} Thu thập các chỉ số hệ thống (CPU, memory, request latency, cache hit rate) từ tất cả các service.
    \item \textbf{Grafana Visualization Dashboard:} Trực quan hóa các chỉ số Prometheus theo thời gian thực, hỗ trợ cài đặt cảnh báo (alert) khi các chỉ số vượt ngưỡng.
\end{itemize}

\subsection{Luồng dữ liệu tổng thể}

Luồng dữ liệu chính của hệ thống có thể tóm tắt qua các bước sau:

\begin{enumerate}
    \item TomTom API và OSM API cung cấp dữ liệu thô cho Data Ingestion Layer.
    \item Pipeline Kappa xử lý, làm sạch, map-match và trích xuất đặc trưng, lưu kết quả vào S3 (.npz) và TimescaleDB.
    \item T-GCN đọc dữ liệu từ Local Disk Cache (.npz), thực hiện huấn luyện và inference, ghi kết quả vào PostgreSQL và Redis.
    \item FastAPI Backend đọc từ Redis (cache hit) hoặc PostgreSQL/TimescaleDB (cache miss), trả về response cho Frontend.
    \item ReactJS Frontend hiển thị bản đồ Folium với lazy loading và overlay tương quan theo yêu cầu người dùng.
\end{enumerate}

\clearpage
\subsection{Thiết kế cơ sở dữ liệu}

Hệ thống UTraffic sử dụng \textbf{PostgreSQL} (kết hợp extension \textbf{PostGIS} và \textbf{TimescaleDB}) làm cơ sở dữ liệu chính. Thiết kế gồm 15 bảng được tổ chức thành 5 nhóm chức năng: xác thực người dùng, topo học đồ thị đường bộ, dữ liệu giao thông thời gian thực, học máy và phiên bản mô hình, phân tích tương quan. Sơ đồ ERD tổng thể được trình bày trong Hình \ref{fig:erd}.

\begin{figure}[H]
    \centering
    \includegraphics[width=\textwidth]{images/erd.png}
    \caption{Sơ đồ ERD cơ sở dữ liệu hệ thống UTraffic}
    \label{fig:erd}
\end{figure}

\subsubsection{Nhóm bảng xác thực người dùng (Authentication)}

\paragraph{Bảng \texttt{users}}
Lưu trữ thông tin tài khoản người dùng hệ thống. Hệ thống hỗ trợ hai vai trò: \texttt{admin} (quản trị viên, có quyền xem toàn bộ dữ liệu phân tích và quản lý mô hình) và \texttt{user} (người dùng thông thường, có quyền xem bản đồ và kết quả tương quan).

\begin{table}[H]
\centering
\caption{Cấu trúc bảng \texttt{users}}
\begin{tabular}{|l|l|l|p{6cm}|}
\hline
\textbf{Cột} & \textbf{Kiểu dữ liệu} & \textbf{Ràng buộc} & \textbf{Mô tả} \\
\hline
\texttt{id} & UUID & PK & Định danh duy nhất \\
\texttt{email} & VARCHAR & UK, NOT NULL & Email đăng nhập \\
\texttt{password\_hash} & VARCHAR & NOT NULL & Mật khẩu băm (bcrypt) \\
\texttt{full\_name} & VARCHAR & & Tên hiển thị \\
\texttt{role} & VARCHAR & NOT NULL & Vai trò: admin hoặc user \\
\texttt{is\_active} & BOOLEAN & DEFAULT TRUE & Trạng thái tài khoản \\
\texttt{created\_at} & TIMESTAMP & NOT NULL & Thời điểm tạo tài khoản \\
\texttt{last\_login\_at} & TIMESTAMP & & Lần đăng nhập cuối \\
\hline
\end{tabular}
\end{table}

\paragraph{Bảng \texttt{refresh\_tokens}}
Quản lý refresh token theo cơ chế \textbf{Refresh Token Rotation}: mỗi lần làm mới access token, refresh token cũ bị thu hồi (\texttt{is\_revoked = true}) và một refresh token mới được cấp phát. Cơ chế này phát hiện tấn công replay attack khi một token cũ bị sử dụng lại.

\begin{table}[H]
\centering
\caption{Cấu trúc bảng \texttt{refresh\_tokens}}
\begin{tabular}{|l|l|l|p{5.5cm}|}
\hline
\textbf{Cột} & \textbf{Kiểu dữ liệu} & \textbf{Ràng buộc} & \textbf{Mô tả} \\
\hline
\texttt{id} & UUID & PK & Định danh duy nhất \\
\texttt{user\_id} & UUID & FK $\to$ users & Liên kết người dùng \\
\texttt{token\_hash} & VARCHAR & NOT NULL & Hash của refresh token \\
\texttt{device\_info} & VARCHAR & & Thông tin thiết bị \\
\texttt{ip\_address} & VARCHAR & & Địa chỉ IP đăng nhập \\
\texttt{is\_revoked} & BOOLEAN & DEFAULT FALSE & Đã thu hồi hay chưa \\
\texttt{expires\_at} & TIMESTAMP & NOT NULL & Thời điểm hết hạn \\
\texttt{created\_at} & TIMESTAMP & NOT NULL & Thời điểm cấp phát \\
\texttt{revoked\_at} & TIMESTAMP & & Thời điểm thu hồi \\
\hline
\end{tabular}
\end{table}

\paragraph{Bảng \texttt{user\_sessions}}
Theo dõi phiên hoạt động của người dùng trên từng thiết bị. Cho phép admin xem danh sách thiết bị đang đăng nhập và thu hồi phiên từ xa khi cần.

\subsubsection{Nhóm bảng topo học đồ thị đường bộ (Road Graph)}

Nhóm này lưu trữ cấu trúc đồ thị đường bộ sau quá trình map matching giữa OSM và TomTom. Đây là nhóm bảng trung tâm, được tham chiếu bởi hầu hết các nhóm còn lại.

\paragraph{Bảng \texttt{nodes}}
Mỗi node tương ứng với một giao lộ (ngã ba, ngã tư, v.v.) trong subgraph đã được map match. Cột \texttt{node\_index} (0--384) là chỉ số số nguyên tương ứng với vị trí của node trong tensor T-GCN, đảm bảo căn chỉnh (alignment) chính xác giữa cơ sở dữ liệu và mô hình học máy.

\begin{table}[H]
\centering
\caption{Cấu trúc bảng \texttt{nodes}}
\begin{tabular}{|l|l|l|p{5cm}|}
\hline
\textbf{Cột} & \textbf{Kiểu dữ liệu} & \textbf{Ràng buộc} & \textbf{Mô tả} \\
\hline
\texttt{id} & UUID & PK & Định danh nội bộ \\
\texttt{osm\_node\_id} & BIGINT & UK, NOT NULL & ID node trong OSM \\
\texttt{node\_index} & INTEGER & NOT NULL & Chỉ số tensor (0--384) \\
\texttt{lat} & FLOAT & NOT NULL & Vĩ độ \\
\texttt{lon} & FLOAT & NOT NULL & Kinh độ \\
\texttt{geom} & GEOMETRY(POINT) & NOT NULL & Điểm địa lý (PostGIS) \\
\texttt{degree} & FLOAT & & Bậc đỉnh trong subgraph \\
\texttt{betweenness\_norm} & FLOAT & & Betweenness centrality \\
\texttt{lat\_norm} & FLOAT & & Vĩ độ chuẩn hóa [0,1] \\
\texttt{lon\_norm} & FLOAT & & Kinh độ chuẩn hóa [0,1] \\
\texttt{street\_name} & VARCHAR & & Tên đường gần nhất \\
\texttt{created\_at} & TIMESTAMP & NOT NULL & Thời điểm tạo bản ghi \\
\hline
\end{tabular}
\end{table}

\textbf{Index:} Bắt buộc tạo GIST index trên cột \texttt{geom} để hỗ trợ truy vấn bounding box từ Folium lazy loading: \texttt{CREATE INDEX ON nodes USING GIST (geom)}.

\paragraph{Bảng \texttt{edges}}
Mỗi edge là một cạnh có hướng trong subgraph, tương ứng với một đoạn đường một chiều trong OSM. Một tuyến đường hai chiều được biểu diễn bằng hai edge ngược chiều nhau.

\begin{table}[H]
\centering
\caption{Cấu trúc bảng \texttt{edges}}
\begin{tabular}{|l|l|l|p{5cm}|}
\hline
\textbf{Cột} & \textbf{Kiểu dữ liệu} & \textbf{Ràng buộc} & \textbf{Mô tả} \\
\hline
\texttt{id} & UUID & PK & Định danh duy nhất \\
\texttt{source\_node\_id} & UUID & FK $\to$ nodes & Node đầu cạnh \\
\texttt{target\_node\_id} & UUID & FK $\to$ nodes & Node cuối cạnh \\
\texttt{source\_index} & INTEGER & NOT NULL & Chỉ số tensor node đầu \\
\texttt{target\_index} & INTEGER & NOT NULL & Chỉ số tensor node cuối \\
\texttt{geom} & GEOMETRY(LINESTRING) & NOT NULL & Hình học cạnh (PostGIS) \\
\texttt{length\_m} & FLOAT & & Chiều dài (mét) \\
\texttt{maxspeed\_kmh} & FLOAT & & Tốc độ giới hạn (km/h) \\
\texttt{lanes} & FLOAT & & Số làn xe \\
\texttt{highway\_type} & INTEGER & & Phân loại đường OSM \\
\texttt{created\_at} & TIMESTAMP & NOT NULL & Thời điểm tạo bản ghi \\
\hline
\end{tabular}
\end{table}

\paragraph{Bảng \texttt{tomtom\_segments} và \texttt{segment\_edge\_mappings}}
Bảng \texttt{tomtom\_segments} lưu thông tin gốc của phân đoạn TomTom (ID, hình học LINESTRING, tên đường, FRC, tốc độ giới hạn, chiều dài). Bảng \texttt{segment\_edge\_mappings} là bảng junction thể hiện mối quan hệ nhiều--nhiều giữa phân đoạn TomTom và cạnh OSM: một phân đoạn TomTom dài có thể được map sang nhiều cạnh OSM, và một cạnh OSM có thể nhận dữ liệu từ nhiều phân đoạn TomTom lân cận. Cột \texttt{match\_dist\_m} ghi khoảng cách snap (mét) để đánh giá chất lượng map matching.

\paragraph{Bảng \texttt{graph\_snapshots}}
Lưu phiên bản hóa cấu trúc đồ thị. Mỗi khi pipeline map matching được chạy lại (thêm dữ liệu mới hoặc cập nhật OSM), một snapshot mới được tạo. Cột \texttt{is\_active} xác định phiên bản đồ thị đang được sử dụng bởi mô hình hiện hành. Cột \texttt{npz\_path} trỏ tới file \texttt{.npz} trên S3 hoặc Local Disk chứa ma trận kề và đặc trưng node.

\subsubsection{Nhóm bảng dữ liệu giao thông (Traffic Data)}

\paragraph{Bảng \texttt{traffic\_observations}}
Đây là bảng lớn nhất và quan trọng nhất về mặt dữ liệu, được cấu hình là \textbf{TimescaleDB hypertable} phân vùng theo cột \texttt{observed\_at} với kích thước chunk 1 tuần. Mỗi bản ghi lưu trạng thái giao thông tại một cạnh trong một khung giờ 15 phút.

\begin{table}[H]
\centering
\caption{Cấu trúc bảng \texttt{traffic\_observations} (TimescaleDB hypertable)}
\begin{tabular}{|l|l|l|p{5cm}|}
\hline
\textbf{Cột} & \textbf{Kiểu dữ liệu} & \textbf{Ràng buộc} & \textbf{Mô tả} \\
\hline
\texttt{id} & UUID & PK & Định danh duy nhất \\
\texttt{edge\_id} & UUID & FK $\to$ edges & Cạnh được quan sát \\
\texttt{observed\_at} & TIMESTAMP & NOT NULL & Thời điểm quan sát \\
\texttt{time\_slot} & VARCHAR & & Tên slot (VD: Slot\_0700) \\
\texttt{average\_speed} & FLOAT & & Tốc độ trung bình (km/h) \\
\texttt{harmonic\_avg\_speed} & FLOAT & & Tốc độ điều hòa (km/h) \\
\texttt{median\_speed} & FLOAT & & Tốc độ trung vị (km/h) \\
\texttt{std\_speed} & FLOAT & & Độ lệch chuẩn tốc độ \\
\texttt{avg\_travel\_time} & FLOAT & & Thời gian di chuyển TB (s) \\
\texttt{travel\_time\_ratio} & FLOAT & & Tỉ lệ thời gian/lý tưởng \\
\texttt{congestion\_index} & FLOAT & & Chỉ số tắc nghẽn [0,1] \\
\texttt{speed\_limit\_ratio} & FLOAT & & Tỉ lệ tốc độ/giới hạn \\
\texttt{sample\_size} & INTEGER & & Số lượng mẫu xe \\
\hline
\end{tabular}
\end{table}

\textbf{Lưu ý triển khai:} Sau khi tạo bảng, cần chuyển sang hypertable với câu lệnh:
\begin{verbatim}
SELECT create_hypertable('traffic_observations', 'observed_at',
    chunk_time_interval => INTERVAL '1 week');
\end{verbatim}

Tốc độ ghi ước tính: với 385 node × 24 khung giờ/ngày ≈ 9.240 bản ghi/ngày (ở level edge sau map matching còn nhiều hơn). TimescaleDB continuous aggregate được định nghĩa cho bảng \texttt{hourly\_traffic\_stats} để tính toán thống kê theo giờ mà không cần query toàn bộ bảng gốc.

\paragraph{Bảng \texttt{hourly\_traffic\_stats}}
Bảng aggregate tổng hợp theo giờ, được duy trì tự động bởi TimescaleDB continuous aggregate policy. Khóa chính composite là \texttt{(bucket, edge\_id)}, trong đó \texttt{bucket} là timestamp đã làm tròn theo giờ bằng hàm \texttt{time\_bucket('1 hour', observed\_at)}. Bảng này bổ sung thêm cột \texttt{node\_id} để hỗ trợ truy vấn thống kê trực tiếp theo node mà không cần join với bảng \texttt{edges}.

\subsubsection{Nhóm bảng học máy và phiên bản mô hình (ML \& Model Versioning)}

\paragraph{Bảng \texttt{model\_versions}}
Lưu thông tin về các phiên bản mô hình đã huấn luyện. Cột \texttt{metrics\_by\_horizon} là JSONB chứa metrics tổng hợp của cả 4 horizon. Cột \texttt{checkpoint\_path} trỏ tới file \texttt{best\_model.pth} trên S3 hoặc Local Disk. Chỉ một bản ghi có \texttt{is\_active = true} được sử dụng cho inference ở bất kỳ thời điểm nào.

\paragraph{Bảng \texttt{model\_horizon\_metrics}}
Tách riêng metrics theo từng horizon (1--4) để hỗ trợ so sánh hiệu năng mô hình theo mốc dự báo cụ thể. Với mỗi \texttt{model\_version\_id}, bảng này có đúng 4 bản ghi tương ứng horizon 1, 2, 3, 4.

\begin{table}[H]
\centering
\caption{Cấu trúc bảng \texttt{model\_horizon\_metrics}}
\begin{tabular}{|l|l|l|p{5cm}|}
\hline
\textbf{Cột} & \textbf{Kiểu dữ liệu} & \textbf{Ràng buộc} & \textbf{Mô tả} \\
\hline
\texttt{id} & UUID & PK & Định danh duy nhất \\
\texttt{model\_version\_id} & UUID & FK $\to$ model\_versions & Phiên bản mô hình \\
\texttt{horizon\_step} & INTEGER & NOT NULL & Bước dự báo (1--4) \\
\texttt{rmse} & FLOAT & & Root Mean Squared Error \\
\texttt{mae} & FLOAT & & Mean Absolute Error \\
\texttt{accuracy} & FLOAT & & Accuracy (%) \\
\texttt{r2} & FLOAT & & Hệ số xác định R\textsuperscript{2} \\
\texttt{var\_score} & FLOAT & & Explained Variance Score \\
\hline
\end{tabular}
\end{table}

Kết quả thực nghiệm từ \texttt{comparison\_results.json} cho thấy T-GCN đạt hiệu năng vượt trội so với các baseline tại tất cả 4 horizon, cụ thể tại horizon 1: RMSE = 0.443, MAE = 0.337, R\textsuperscript{2} = 0.797; và tại horizon 4: RMSE = 0.560, MAE = 0.406, R\textsuperscript{2} = 0.725.

\paragraph{Bảng \texttt{predictions}}
Lưu kết quả dự báo tốc độ từ T-GCN. Mỗi lần chạy prediction job sinh ra $385 \times 4 = 1.540$ bản ghi (385 node × 4 horizon). Với tần suất 15 phút/lần, bảng này tích lũy khoảng 147.840 bản ghi/ngày. Để kiểm soát kích thước, bảng này cần được cấu hình là \textbf{TimescaleDB hypertable} với TTL policy tự động xóa dữ liệu cũ hơn 30 ngày.

\begin{table}[H]
\centering
\caption{Cấu trúc bảng \texttt{predictions}}
\begin{tabular}{|l|l|l|p{5cm}|}
\hline
\textbf{Cột} & \textbf{Kiểu dữ liệu} & \textbf{Ràng buộc} & \textbf{Mô tả} \\
\hline
\texttt{id} & UUID & PK & Định danh duy nhất \\
\texttt{model\_version\_id} & UUID & FK $\to$ model\_versions & Mô hình sinh dự báo \\
\texttt{node\_id} & UUID & FK $\to$ nodes & Node được dự báo \\
\texttt{predicted\_at} & TIMESTAMP & NOT NULL & Thời điểm chạy inference \\
\texttt{target\_time} & TIMESTAMP & NOT NULL & Thời điểm dự báo tới \\
\texttt{horizon\_step} & INTEGER & NOT NULL & Bước dự báo (1--4) \\
\texttt{predicted\_speed} & FLOAT & & Tốc độ dự báo (km/h) \\
\texttt{congestion\_index} & FLOAT & & Chỉ số tắc nghẽn dự báo \\
\texttt{congestion\_label} & VARCHAR & & Nhãn: free/slow/congested \\
\hline
\end{tabular}
\end{table}

\textbf{Index:} Cần tạo composite index \texttt{(node\_id, predicted\_at, horizon\_step)} để tối ưu truy vấn từ Prediction Service.

\subsubsection{Nhóm bảng phân tích tương quan (Correlation Analysis)}

Đây là nhóm bảng phục vụ tính năng cốt lõi của hệ thống: hiển thị mối tương quan giữa các nút giao thông khi người dùng click vào bản đồ.

\paragraph{Bảng \texttt{correlation\_snapshots}}
Lưu metadata của mỗi lần chạy correlation job. Cột \texttt{npz\_path} trỏ tới file ma trận tương quan đầy đủ $385 \times 385$ trên Local Disk/S3. Cột \texttt{threshold\_stored} ghi nhận ngưỡng lọc đã áp dụng khi ghi vào bảng \texttt{node\_correlations} (mặc định $\geq 0.7$). Chỉ một snapshot có \texttt{is\_active = true} được serve bởi API tại bất kỳ thời điểm nào.

\paragraph{Bảng \texttt{node\_correlations}}
Lưu các cặp node có mức độ tương quan đáng kể. Từ dữ liệu thực tế, với ngưỡng $|corr| \geq 0.7$, bảng này chứa khoảng 34.307 cặp. Hai cột \texttt{rank\_from\_a} và \texttt{rank\_from\_b} ghi nhận thứ hạng của cặp này trong danh sách láng giềng của node A và node B tương ứng (xếp theo chiều giảm dần của $|corr|$), hỗ trợ API trả về top-K neighbors một cách hiệu quả.

\begin{table}[H]
\centering
\caption{Cấu trúc bảng \texttt{node\_correlations}}
\begin{tabular}{|l|l|l|p{5cm}|}
\hline
\textbf{Cột} & \textbf{Kiểu dữ liệu} & \textbf{Ràng buộc} & \textbf{Mô tả} \\
\hline
\texttt{id} & UUID & PK & Định danh duy nhất \\
\texttt{snapshot\_id} & UUID & FK $\to$ correlation\_snapshots & Phiên bản correlation \\
\texttt{node\_a\_id} & UUID & FK $\to$ nodes & Node thứ nhất \\
\texttt{node\_b\_id} & UUID & FK $\to$ nodes & Node thứ hai \\
\texttt{correlation\_value} & FLOAT & NOT NULL & Hệ số tương quan [-1,1] \\
\texttt{rank\_from\_a} & INTEGER & & Thứ hạng từ góc nhìn A \\
\texttt{rank\_from\_b} & INTEGER & & Thứ hạng từ góc nhìn B \\
\texttt{is\_adjacent} & BOOLEAN & DEFAULT FALSE & Có cạnh trực tiếp không \\
\texttt{created\_at} & TIMESTAMP & NOT NULL & Thời điểm tạo bản ghi \\
\hline
\end{tabular}
\end{table}

\textbf{Index cần thiết:}
\begin{verbatim}
CREATE INDEX ON node_correlations (node_a_id, correlation_value DESC);
CREATE INDEX ON node_correlations (node_b_id, correlation_value DESC);
\end{verbatim}

\paragraph{Bảng \texttt{node\_correlation\_cache}}
Đây là bảng tối ưu hóa hiệu năng quan trọng nhất cho tính năng click-node. Thay vì query và join bảng \texttt{node\_correlations} mỗi lần người dùng click, hệ thống pre-aggregate sẵn top-20 láng giềng của mỗi node vào một JSONB blob và cache trong bảng này. Mỗi blob có dạng:

\begin{verbatim}
[
  {"node_id": "uuid-...", "corr": 0.9979,
   "lat": 10.7650, "lon": 106.6781, "rank": 1},
  {"node_id": "uuid-...", "corr": 0.9975,
   "lat": 10.7623, "lon": 106.6754, "rank": 2},
  ...
]
\end{verbatim}

Tổng kích thước cache cho 385 node ước tính $\sim$0.63 MB, đủ nhỏ để tải toàn bộ vào Redis-Corr khi khởi động. Khi người dùng click vào node X, Correlation Service chỉ cần một lệnh \texttt{GET} từ Redis, thời gian phản hồi dưới 1ms.

\subsubsection{Chiến lược index và tối ưu hóa}

Ngoài các index đã đề cập, hệ thống áp dụng các chiến lược tối ưu sau:

\begin{itemize}
    \item \textbf{Partial index cho correlation:} Tạo partial index trên bảng \texttt{node\_correlations} chỉ cho các bản ghi thuộc snapshot đang active, giảm kích thước index khi có nhiều snapshot lịch sử.
    \item \textbf{BRIN index cho time-series:} Bảng \texttt{traffic\_observations} sử dụng BRIN index trên cột \texttt{observed\_at} thay vì B-tree do dữ liệu được ghi tuần tự theo thời gian, BRIN có kích thước nhỏ hơn đáng kể.
    \item \textbf{Connection pooling:} Sử dụng \textbf{PgBouncer} để quản lý connection pool giữa FastAPI services và PostgreSQL, tránh tình trạng connection exhaustion khi nhiều request đồng thời.
    \item \textbf{Read replica:} Trong môi trường production, TimescaleDB replica được cấu hình để tách tải đọc (analytics queries từ Dashboard) khỏi tải ghi (pipeline sync), đảm bảo hiệu năng cho cả hai luồng.
\end{itemize}

\subsubsection{Tóm tắt thiết kế cơ sở dữ liệu}

\begin{table}[H]
\centering
\caption{Tổng hợp các bảng trong hệ thống UTraffic}
\begin{tabular}{|l|l|p{6cm}|}
\hline
\textbf{Nhóm} & \textbf{Bảng} & \textbf{Chức năng chính} \\
\hline
\multirow{3}{*}{Authentication} & \texttt{users} & Tài khoản người dùng \\
 & \texttt{refresh\_tokens} & Quản lý refresh token \\
 & \texttt{user\_sessions} & Theo dõi phiên đăng nhập \\
\hline
\multirow{4}{*}{Road Graph} & \texttt{nodes} & Node giao lộ (OSM + PostGIS) \\
 & \texttt{edges} & Cạnh đường bộ có hướng \\
 & \texttt{tomtom\_segments} & Phân đoạn TomTom gốc \\
 & \texttt{segment\_edge\_mappings} & Kết quả map matching \\
 & \texttt{graph\_snapshots} & Phiên bản hóa đồ thị \\
\hline
\multirow{2}{*}{Traffic Data} & \texttt{traffic\_observations} & Quan sát giao thông (hypertable) \\
 & \texttt{hourly\_traffic\_stats} & Thống kê aggregate theo giờ \\
\hline
\multirow{3}{*}{ML \& Models} & \texttt{model\_versions} & Phiên bản mô hình T-GCN \\
 & \texttt{model\_horizon\_metrics} & Metrics theo từng horizon \\
 & \texttt{predictions} & Dự báo tốc độ (hypertable) \\
\hline
\multirow{3}{*}{Correlation} & \texttt{correlation\_snapshots} & Metadata phiên phân tích \\
 & \texttt{node\_correlations} & Ma trận tương quan (lọc) \\
 & \texttt{node\_correlation\_cache} & Cache top-20 JSONB per node \\
\hline
\end{tabular}
\end{table}